\subsection{Annihilators and Transposes}
Throughout this subsection, all vector spaces are over the same fixed field $\mathbb{F}$.

\begin{definition} \label{def:annihilator} \hlt{Annihilator}\\
If $\mathcal{U}\subseteq\mathcal{V}$, the \hlt{annihilator of $\mathcal{U}$ in $\mathcal{V}'$} is $\mathcal{U}^0\defeq\{\varphi\in\mathcal{V}':(\forall\,\mathbf{u}\in\mathcal{U})\,\varphi(\mathbf{u})=0\}$.\\
The ambient space $\mathcal{V}$ is understood from context; changing it changes the dual space containing $\mathcal{U}^0$.
\end{definition}

\begin{proposition} \label{prop:annihilator-is-subspace} \hlt{Annihilators Are Subspaces}\\
If $\mathcal{U}\subseteq\mathcal{V}$, then $\mathcal{U}^0$ is a subspace of $\mathcal{V}'$.
\begin{proof}
The zero functional belongs to $\mathcal{U}^0$. If $\varphi,\psi\in\mathcal{U}^0$, $a\in\mathbb{F}$, and $\mathbf{u}\in\mathcal{U}$, then $(a\varphi+\psi)(\mathbf{u})=a\varphi(\mathbf{u})+\psi(\mathbf{u})=0$.\\
Hence $a\varphi+\psi\in\mathcal{U}^0$, and $\mathcal{U}^0\leq\mathcal{V}'$ by Theorem \ref{thm:subspacetest}.\\
\hlt{Structural Commentary}: The subspace test reduces closure to the preservation of pointwise vanishing.\\
\hlt{Key Insight}: Linear combinations preserve every annihilating constraint.
\end{proof}
\end{proposition}

\begin{theorem} \label{thm:dimension-of-annihilator} \hlt{Dimension of the Annihilator}\\
Let $\mathcal{V}$ be finite-dimensional and let $\mathcal{U}$ be a subspace of $\mathcal{V}$. Then $\dim(\mathcal{U})+\dim(\mathcal{U}^0)=\dim(\mathcal{V})$.
\begin{proof}
Choose basis $(\mathbf{u}_1,\ldots,\mathbf{u}_m)$ of $\mathcal{U}$ by Corollary \ref{cor:finite-dimensional-spaces-and-subspaces-have-bases}, extend to basis $\beta=(\mathbf{u}_1,\ldots,\mathbf{u}_m,\mathbf{w}_1,\ldots,\mathbf{w}_k)$ of $\mathcal{V}$ by Theorem \ref{thm:extending-independent-lists}.\\
Let $\beta'=(\varphi_1,\ldots,\varphi_m,\psi_1,\ldots,\psi_k)$ be the dual basis.\\
By Theorem \ref{thm:dual-basis-and-dimension}(i) and Definition \ref{def:dual-basis}, every $\varphi\in\mathcal{V}'$ satisfies $\varphi=\sum\limits_{j=1}^{m}\varphi(\mathbf{u}_j)\varphi_j+\sum\limits_{\ell=1}^{k}\varphi(\mathbf{w}_\ell)\psi_\ell$. Thus $\varphi\in\mathcal{U}^0$ implies $\varphi=\sum\limits_{\ell=1}^{k}\varphi(\mathbf{w}_\ell)\psi_\ell$.\\
Conversely, each $\psi_\ell$ vanishes on $\mathcal{U}=\operatorname{span}(\mathbf{u}_1,\ldots,\mathbf{u}_m)$ by Definition \ref{def:dual-basis}.\\
Hence $(\psi_1,\ldots,\psi_k)$ spans $\mathcal{U}^0$ and is linearly independent as a sublist of $\beta'$, so it is a basis of $\mathcal{U}^0$.\\
The adapted and annihilator bases have lengths $m+k$ and $k$, respectively, so $\dim(\mathcal{U})+\dim(\mathcal{U}^0)=m+k=\dim(\mathcal{V})$.\\
\hlt{Structural Commentary}: An adapted basis identifies the annihilator with the dual coordinates outside $\mathcal{U}$.\\
\hlt{Key Insight}: Vanishing deletes precisely the dual coordinates indexed by $\mathcal{U}$.
\end{proof}
\end{theorem}

\begin{theorem} \label{thm:annihilator-lattice-identities} \hlt{Annihilator Lattice Identities}\\
Let $\mathcal{V}$ be finite-dimensional, and let $\mathcal{S},\mathcal{T}$ be subspaces of $\mathcal{V}$.
\begin{enumerate}[label=(\roman*)]
\setlength{\itemsep}{0pt}
\item $\mathcal{S}\subseteq\mathcal{T}\Rightarrow\mathcal{T}^0\subseteq\mathcal{S}^0$.
\item $(\mathcal{S}+\mathcal{T})^0=\mathcal{S}^0\cap\mathcal{T}^0$.
\item $(\mathcal{S}\cap\mathcal{T})^0=\mathcal{S}^0+\mathcal{T}^0$.
\end{enumerate}
\begin{proof}
\begin{enumerate}[label=(\roman*)]
\setlength{\itemsep}{0pt}
\item If $\mathcal{S}\subseteq\mathcal{T}$ and $\varphi\in\mathcal{T}^0$, then $(\forall\,\mathbf{s}\in\mathcal{S})\,\varphi(\mathbf{s})=0$, so $\varphi\in\mathcal{S}^0$.
\item If $\varphi\in(\mathcal{S}+\mathcal{T})^0$, then $\varphi|_{\mathcal{S}}=0$ and $\varphi|_{\mathcal{T}}=0$, so $\varphi\in\mathcal{S}^0\cap\mathcal{T}^0$.\\
Conversely, if $\varphi\in\mathcal{S}^0\cap\mathcal{T}^0$, then $(\forall\,\mathbf{s}\in\mathcal{S})(\forall\,\mathbf{t}\in\mathcal{T})\,\varphi(\mathbf{s}+\mathbf{t})=\varphi(\mathbf{s})+\varphi(\mathbf{t})=0$, so $\varphi\in(\mathcal{S}+\mathcal{T})^0$.
\item By (i), $\mathcal{S}^0+\mathcal{T}^0\subseteq(\mathcal{S}\cap\mathcal{T})^0$. By (ii) and Theorem \ref{thm:dimension-formula-for-sum}, $\dim(\mathcal{S}^0+\mathcal{T}^0)=\dim(\mathcal{S}^0)+\dim(\mathcal{T}^0)-\dim((\mathcal{S}+\mathcal{T})^0)$.\\
The identity $\dim(\mathcal{S}^0+\mathcal{T}^0)=\dim(\mathcal{V})-\dim(\mathcal{S}\cap\mathcal{T})=\dim((\mathcal{S}\cap\mathcal{T})^0)$ holds by Theorem \ref{thm:dimension-of-annihilator} and Theorem \ref{thm:dimension-formula-for-sum}.\\
Hence the inclusion is an equality by Proposition \ref{prop:dimension-bounds-and-basis-criteria}.
\end{enumerate}
\hlt{Structural Commentary}: Two direct inclusions establish order reversal and sums; a dimension comparison recovers intersections.\\
\hlt{Key Insight}: Annihilation reverses the finite-dimensional subspace lattice.
\end{proof}
\end{theorem}

\begin{theorem} \label{thm:double-annihilator-finite-dimensional} \hlt{Double Annihilator}\\
Let $\dim(\mathcal{V})<\infty$ and $M\subseteq\mathcal{V}$, and set $M^{00}\defeq(M^0)^0\subseteq\mathcal{V}''$. Then:
\begin{enumerate}[label=(\roman*)]
\setlength{\itemsep}{0pt}
\item $M^{00}=\Lambda_{\mathcal{V}}(\operatorname{span}(M))$;
\item if $\mathcal{S}\leq\mathcal{V}$, then $\mathcal{S}^{00}=\Lambda_{\mathcal{V}}(\mathcal{S})$, and the corestriction $\Lambda_{\mathcal{V}}|_{\mathcal{S}}:\mathcal{S}\to\mathcal{S}^{00}$ is an isomorphism.
\end{enumerate}
\begin{proof}
\begin{enumerate}[label=(\roman*)]
\setlength{\itemsep}{0pt}
\item If $\mathbf{v}\in\operatorname{span}(M)$ and $\varphi\in M^0$, then $\Lambda_{\mathcal{V}}(\mathbf{v})(\varphi)=\varphi(\mathbf{v})=0$ by linearity of $\varphi$, so $\Lambda_{\mathcal{V}}(\operatorname{span}(M))\subseteq M^{00}$.\\
For the reverse inclusion, let $\eta\in M^{00}$. The relation $(\exists!\,\mathbf{v}\in\mathcal{V})\,\eta=\Lambda_{\mathcal{V}}(\mathbf{v})$ holds by Theorem \ref{thm:algebraic-reflexivity-natural-map}(ii); fix this $\mathbf{v}$.\\
Suppose $\mathbf{v}\notin\operatorname{span}(M)$. Choose a basis $(\mathbf{u}_1,\ldots,\mathbf{u}_m)$ of $\operatorname{span}(M)$ by Corollary \ref{cor:finite-dimensional-spaces-and-subspaces-have-bases}.\\
The list $(\mathbf{u}_1,\ldots,\mathbf{u}_m,\mathbf{v})$ is linearly independent by Proposition \ref{prop:adjoiningandremovingvectors} and extends to a basis of $\mathcal{V}$ by Theorem \ref{thm:extending-independent-lists}.\\
Its dual-basis functional $\varphi$ corresponding to $\mathbf{v}$ satisfies $\varphi|_{\operatorname{span}(M)}=0$ and $\varphi(\mathbf{v})=1$.\\
Thus $\varphi\in M^0$ but $\eta(\varphi)=1$, contrary to $\eta\in M^{00}$. Hence $M^{00}\subseteq\Lambda_{\mathcal{V}}(\operatorname{span}(M))$.
\item The identity $\mathcal{S}^{00}=\Lambda_{\mathcal{V}}(\mathcal{S})$ holds by (i) with $M=\mathcal{S}$. The corestriction is linear and injective by Theorem \ref{thm:algebraic-reflexivity-natural-map}(i)--(ii), and it is surjective onto $\mathcal{S}^{00}$. Therefore it is an isomorphism by Proposition \ref{prop:linear-map-invertible-iff-bijective}.
\end{enumerate}
\hlt{Structural Commentary}: One inclusion extends vanishing; an adapted dual basis separates every vector outside the span.\\
\hlt{Key Insight}: Double annihilation remembers exactly the generated subspace.
\end{proof}
\end{theorem}

\begin{theorem} \label{thm:direct-sum-dual-annihilator} \hlt{Direct Sums and Annihilators}\\
Suppose $\mathcal{V}=\mathcal{S}\oplus\mathcal{T}$, and define $E:\mathcal{T}'\to\mathcal{S}^0$ by $(\forall\,\varphi\in\mathcal{T}')(\forall\,\mathbf{s}\in\mathcal{S})(\forall\,\mathbf{t}\in\mathcal{T})\,(E\varphi)(\mathbf{s}+\mathbf{t})\defeq\varphi(\mathbf{t})$. Then:
\begin{enumerate}[label=(\roman*)]
\setlength{\itemsep}{0pt}
\item $E$ is an isomorphism;
\item $\mathcal{V}'=\mathcal{S}^0\oplus\mathcal{T}^0$.
\end{enumerate}
\begin{proof}
\begin{enumerate}[label=(\roman*)]
\setlength{\itemsep}{0pt}
\item Each $\mathbf{v}\in\mathcal{V}$ has a unique representation $\mathbf{v}=\mathbf{s}+\mathbf{t}$ by Proposition \ref{prop:uniqueness-direct-sum-two-subspaces}, so $E$ is well-defined.\\
If $\mathbf{v}_i=\mathbf{s}_i+\mathbf{t}_i$ for $i=1,2$ and $a,b\in\mathbb{F}$, then $(E\varphi)(a\mathbf{v}_1+b\mathbf{v}_2)=\varphi(a\mathbf{t}_1+b\mathbf{t}_2)=a(E\varphi)(\mathbf{v}_1)+b(E\varphi)(\mathbf{v}_2)$, so $E\varphi\in\mathcal{V}'$.\\
The relation $(\forall\,\mathbf{s}\in\mathcal{S})\,(E\varphi)(\mathbf{s})=0$ places $E\varphi$ in $\mathcal{S}^0$. For $a,b\in\mathbb{F}$ and $\varphi,\psi\in\mathcal{T}'$, the identity $E(a\varphi+b\psi)=aE\varphi+bE\psi$ holds pointwise and hence by Assumption \ref{assump:function-extensionality}, so $E\in\LL(\mathcal{T}',\mathcal{S}^0)$. Define $R:\mathcal{S}^0\to\mathcal{T}'$ by $(\forall\,\eta\in\mathcal{S}^0)\,R\eta\defeq\eta|_{\mathcal{T}}$.\\
If $\eta\in\mathcal{S}^0$, $a,b\in\mathbb{F}$, and $\mathbf{t}_1,\mathbf{t}_2\in\mathcal{T}$, then $(R\eta)(a\mathbf{t}_1+b\mathbf{t}_2)=a(R\eta)(\mathbf{t}_1)+b(R\eta)(\mathbf{t}_2)$ by linearity of $\eta$, so $R\eta\in\mathcal{T}'$.\\
For $a,b\in\mathbb{F}$ and $\eta,\zeta\in\mathcal{S}^0$, the identity $R(a\eta+b\zeta)=aR\eta+bR\zeta$ holds pointwise, so $R\in\LL(\mathcal{S}^0,\mathcal{T}')$ by Assumption \ref{assump:function-extensionality}.\\
The identities $(\forall\,\varphi\in\mathcal{T}')\,RE\varphi=\varphi$ and $(\forall\,\eta\in\mathcal{S}^0)(\forall\,\mathbf{s}\in\mathcal{S})(\forall\,\mathbf{t}\in\mathcal{T})\,(ER\eta)(\mathbf{s}+\mathbf{t})=\eta(\mathbf{t})=\eta(\mathbf{s}+\mathbf{t})$ hold by the definitions of $E$, $R$, and $\mathcal{S}^0$. Thus $R$ is a two-sided inverse of $E$ by Assumption \ref{assump:function-extensionality}, and $E$ is an isomorphism.
\item For $\varphi\in\mathcal{V}'$, set $\varphi_{\mathcal{S}}\defeq\varphi\rho_{\mathcal{S},\mathcal{T}}$ and $\varphi_{\mathcal{T}}\defeq\varphi\rho_{\mathcal{T},\mathcal{S}}$. These functionals belong to $\mathcal{V}'$ by Proposition \ref{prop:algebraic-projection-range-kernel-idempotent}(ii) and Proposition \ref{prop:linear-map-product-algebra}(i).\\
Their respective membership in $\mathcal{T}^0$ and $\mathcal{S}^0$ holds by Proposition \ref{prop:algebraic-projection-range-kernel-idempotent}(iii).\\
The identity $\varphi=\varphi_{\mathcal{S}}+\varphi_{\mathcal{T}}$ holds by Proposition \ref{prop:complementary-algebraic-projections}(i), so $\mathcal{V}'=\mathcal{S}^0+\mathcal{T}^0$.\\
If $\eta\in\mathcal{S}^0\cap\mathcal{T}^0$, then $(\forall\,\mathbf{s}\in\mathcal{S})(\forall\,\mathbf{t}\in\mathcal{T})\,\eta(\mathbf{s}+\mathbf{t})=0$, so $\eta=0_{\mathcal{V}'}$ by Assumption \ref{assump:function-extensionality}. Thus $\mathcal{V}'=\mathcal{S}^0\oplus\mathcal{T}^0$.
\end{enumerate}
\hlt{Structural Commentary}: Extension and restriction identify one annihilator; component functionals split every element of the dual.\\
\hlt{Key Insight}: Each dual component ignores the opposite primal summand.
\end{proof}
\end{theorem}

\begin{theorem} \label{thm:dual-kernel-range-annihilator} \hlt{Kernel and Range of the Dual Map}\\
Let $T\in\LL(\mathcal{V},\mathcal{W})$. The following statements hold, with finite-dimensionality of $\mathcal{V}$ and $\mathcal{W}$ required only in (ii)--(iv):
\begin{enumerate}[label=(\roman*)]
\setlength{\itemsep}{0pt}
\item $\ker(T')=(\operatorname{range}(T))^0$;
\item $\dim(\ker(T'))=\dim(\ker(T))+\dim(\mathcal{W})-\dim(\mathcal{V})$;
\item $\operatorname{range}(T')=(\ker(T))^0$;
\item $\dim(\operatorname{range}(T'))=\dim(\operatorname{range}(T))$.
\end{enumerate}
\begin{proof}
For (ii)--(iv), assume that $\mathcal{V}$ and $\mathcal{W}$ are finite-dimensional.
\begin{enumerate}[label=(\roman*)]
\setlength{\itemsep}{0pt}
\item For $\varphi\in\mathcal{W}'$, $T'\varphi=0$ iff $(\forall\,\mathbf{v}\in\mathcal{V})\,\varphi(T\mathbf{v})=0$ iff $\varphi\in(\operatorname{range}(T))^0$.
\item By (i), Theorem \ref{thm:dimension-of-annihilator}, and Theorem \ref{thm:rank-nullity}(ii), $\dim(\ker(T'))=\dim(\mathcal{W})-\dim(\operatorname{range}(T))=\dim(\ker(T))+\dim(\mathcal{W})-\dim(\mathcal{V})$.
\item If $\varphi=T'\psi$ and $\mathbf{v}\in\ker(T)$, then $\varphi(\mathbf{v})=\psi(T\mathbf{v})=0$, so $\operatorname{range}(T')\subseteq(\ker(T))^0$.\\
By Theorem \ref{thm:rank-nullity}(ii), (ii), and Theorem \ref{thm:dual-basis-and-dimension}(ii), $\dim(\operatorname{range}(T'))=\dim(\mathcal{W}')-\dim(\ker(T'))=\dim(\mathcal{V})-\dim(\ker(T))$.\\
The latter dimension equals $\dim((\ker(T))^0)$ by Theorem \ref{thm:dimension-of-annihilator}, so the inclusion is an equality by Proposition \ref{prop:dimension-bounds-and-basis-criteria}.
\item By (iii), Theorem \ref{thm:dimension-of-annihilator}, and Theorem \ref{thm:rank-nullity}(ii), $\dim(\operatorname{range}(T'))=\dim(\mathcal{V})-\dim(\ker(T))=\dim(\operatorname{range}(T))$.
\end{enumerate}
\hlt{Structural Commentary}: A direct kernel identity and finite-dimensional counts determine the dual range and both dimensions.\\
\hlt{Key Insight}: Pullback exchanges unreachable vectors with annihilating functionals.
\end{proof}
\end{theorem}

\begin{corollary} \label{cor:surjective-iff-dual-injective} \hlt{Surjectivity and Dual Injectivity}\\
Let $\mathcal{V}$ and $\mathcal{W}$ be finite-dimensional, and let $T\in\LL(\mathcal{V},\mathcal{W})$. Then $T$ is surjective iff $T'$ is injective.
\begin{proof}
\begin{enumerate}[label=\roman*.]
\setlength{\itemsep}{0pt}
\item[$(\Rightarrow)$] If $T$ is surjective, then $\operatorname{range}(T)=\mathcal{W}$, so $\ker(T')=\mathcal{W}^0=\{0_{\mathcal{W}'}\}$ by Theorem \ref{thm:dual-kernel-range-annihilator}(i). Hence $T'$ is injective by Proposition \ref{prop:injective-iff-kernel-zero}.
\item[$(\Leftarrow)$] If $T'$ is injective, then $(\operatorname{range}(T))^0=\ker(T')=\{0_{\mathcal{W}'}\}$ by Proposition \ref{prop:injective-iff-kernel-zero} and Theorem \ref{thm:dual-kernel-range-annihilator}(i). \\
Thus $\dim(\operatorname{range}(T))=\dim(\mathcal{W})$ by Theorem \ref{thm:dimension-of-annihilator}, so $\operatorname{range}(T)=\mathcal{W}$ by Proposition \ref{prop:dimension-bounds-and-basis-criteria}.
\end{enumerate}
\hlt{Structural Commentary}: The dual-kernel identity translates surjectivity into a zero-kernel criterion.\\
\hlt{Key Insight}: A full range has no nonzero annihilator.
\end{proof}
\end{corollary}

\begin{corollary} \label{cor:injective-iff-dual-surjective} \hlt{Injectivity and Dual Surjectivity}\\
Let $\mathcal{V}$ and $\mathcal{W}$ be finite-dimensional, and let $T\in\LL(\mathcal{V},\mathcal{W})$. Then $T$ is injective iff $T'$ is surjective.
\begin{proof}
\begin{enumerate}[label=\roman*.]
\setlength{\itemsep}{0pt}
\item[$(\Rightarrow)$] If $T$ is injective, then $T'$ is surjective by Proposition \ref{prop:injective-iff-kernel-zero}, Theorem \ref{thm:dual-kernel-range-annihilator}(iii), and Definition \ref{def:annihilator}.
\item[$(\Leftarrow)$] If $T'$ is surjective, then $(\ker(T))^0=\operatorname{range}(T')=\mathcal{V}'$ by Theorem \ref{thm:dual-kernel-range-annihilator}(iii).\\
Thus $\dim(\ker(T))=0$ by Theorem \ref{thm:dimension-of-annihilator}, so $\ker(T)=\{\mathbf{0}_{\mathcal{V}}\}$ by Proposition \ref{prop:dimension-bounds-and-basis-criteria}. Therefore $T$ is injective by Proposition \ref{prop:injective-iff-kernel-zero}.
\end{enumerate}
\hlt{Structural Commentary}: The dual-range identity converts injectivity into surjectivity of the pullback.\\
\hlt{Key Insight}: A trivial kernel leaves no functional-extension obstruction.
\end{proof}
\end{corollary}

\begin{corollary} \label{cor:quotient-dual-annihilator-isomorphism} \hlt{Dual of a Quotient as an Annihilator}\\
Let $\mathcal{V}$ be finite-dimensional, let $\mathcal{U}\leq\mathcal{V}$, and let $\pi:\mathcal{V}\to\mathcal{V}/\mathcal{U}$ be the quotient map. Then:
\begin{enumerate}[label=(\roman*)]
\setlength{\itemsep}{0pt}
\item $\ker(\pi')=\{0_{(\mathcal{V}/\mathcal{U})'}\}$;
\item $\operatorname{range}(\pi')=\mathcal{U}^0$;
\item the corestriction $\pi':(\mathcal{V}/\mathcal{U})'\to\mathcal{U}^0$ is an isomorphism, so $(\mathcal{V}/\mathcal{U})'\cong\mathcal{U}^0$.
\end{enumerate}
\begin{proof}
The quotient is finite-dimensional by Theorem \ref{thm:dimension-of-quotient}.
\begin{enumerate}[label=(\roman*)]
\setlength{\itemsep}{0pt}
\item The map $\pi$ is surjective, so $\pi'$ is injective by Corollary \ref{cor:surjective-iff-dual-injective}. Hence $\ker(\pi')=\{0_{(\mathcal{V}/\mathcal{U})'}\}$ by Proposition \ref{prop:injective-iff-kernel-zero}.
\item Since $\ker(\pi)=\mathcal{U}$, $\operatorname{range}(\pi')=(\ker(\pi))^0=\mathcal{U}^0$ by Theorem \ref{thm:dual-kernel-range-annihilator}(iii).
\item By (i)--(ii), the corestriction is a linear bijection and hence an isomorphism by Proposition \ref{prop:linear-map-invertible-iff-bijective}.
\end{enumerate}
\hlt{Structural Commentary}: Dual kernel and range identities make the quotient pullback a bijection onto the annihilator.\\
\hlt{Key Insight}: Quotient functionals are exactly functionals vanishing on the collapsed subspace.
\end{proof}
\end{corollary}

\begin{theorem} \label{thm:range-membership-dual-alternative} \hlt{Range Membership Dual Alternative}\\
Let $\mathcal{V},\mathcal{W}$ be finite-dimensional, $T\in\LL(\mathcal{V},\mathcal{W})$, and $\mathbf{w}\in\mathcal{W}$. Exactly one holds:
\begin{enumerate}[label=(\roman*)]
\setlength{\itemsep}{0pt}
\item $\mathbf{w}\in\operatorname{range}(T)$;
\item $(\exists\,\psi\in\mathcal{W}')[\psi T=0\land\psi(\mathbf{w})\neq0]$.
\end{enumerate}
\begin{proof}
\begin{enumerate}[label=(\roman*)]
\setlength{\itemsep}{0pt}
\item If (i) and (ii) hold, choose $\mathbf{v}\in\mathcal{V}$ with $\mathbf{w}=T\mathbf{v}$ and $\psi\in\mathcal{W}'$ satisfying (ii).\\
Then $\psi(\mathbf{w})=\psi(T\mathbf{v})=0$, contrary to $\psi(\mathbf{w})\neq0$. Hence (i) and (ii) cannot both hold.
\item Suppose (i) fails. The map $\Lambda_{\mathcal{W}}$ is injective by Theorem \ref{thm:algebraic-reflexivity-natural-map}(ii), so $\Lambda_{\mathcal{W}}(\mathbf{w})\notin\Lambda_{\mathcal{W}}(\operatorname{range}(T))=(\operatorname{range}(T))^{00}$ by Theorem \ref{thm:double-annihilator-finite-dimensional}(i).\\
Hence $(\exists\,\psi\in(\operatorname{range}(T))^0)\,\Lambda_{\mathcal{W}}(\mathbf{w})(\psi)\neq0$ by Definition \ref{def:annihilator}. For such $\psi$, one has $\psi(\mathbf{w})\neq0$ and $(\forall\,\mathbf{v}\in\mathcal{V})\,(\psi T)(\mathbf{v})=\psi(T\mathbf{v})=0$, so $\psi T=0$ by Assumption \ref{assump:function-extensionality}. Thus (ii) holds whenever (i) fails.
\end{enumerate}
\hlt{Structural Commentary}: Direct evaluation proves exclusion; double annihilation constructs every missing-range certificate.\\
\hlt{Key Insight}: A separating functional certifies every missing range vector.
\end{proof}
\end{theorem}

\begin{illustrationfigure}[H]
  \centering
  \begin{tikzpicture}[
  >=Stealth,
  every node/.style={font=\small},
  ambient/.style={draw=gray!45, fill=gray!3, rounded corners=3pt, line width=0.55pt,
    minimum width=6.5cm, minimum height=3.35cm},
  subspace/.style={draw=teal!60!black, fill=teal!8, rounded corners=2pt,
    line width=0.55pt, minimum width=2.55cm, minimum height=0.72cm, align=center},
  point/.style={circle, fill=blue!70!black, inner sep=1.5pt},
  cert/.style={draw=red!60!black, dashed, line width=0.7pt},
  detect/.style={->, red!60!black, line width=0.7pt}
]
  \node[ambient] at (0,0) {};
  \node[anchor=north west] at (-3.02,1.52) {$\mathcal{W}$};

  \draw[cert] (-2.65,-0.40) -- (2.75,-0.40);
  \node[red!60!black, anchor=south west] at (-2.82,-0.34) {$\ker\psi$};

  \node[subspace] at (-0.55,-0.40) {$\operatorname{range}(T)$\\[-1pt]{\footnotesize $\psi=0$}};

  \node[point] (w) at (1.95,0.92) {};
  \node[anchor=west] at (2.08,0.92) {$\mathbf{w}$};
  \draw[detect] (w) -- (1.95,-0.32);
  \node[red!60!black, anchor=east] at (1.72,0.30) {$\psi(\mathbf{w})\neq0$};
\end{tikzpicture}

  \caption{The dual certificate $\psi$ vanishes on $\operatorname{range}(T)$ but detects $\mathbf{w}$ outside that range.}
  \label{fig:range-membership-dual-certificate}
\end{illustrationfigure}

Figure \ref{fig:range-membership-dual-certificate} illustrates Theorem \ref{thm:range-membership-dual-alternative}: $\operatorname{range}(T)\subseteq\ker(\psi)$ while $\mathbf{w}\notin\ker(\psi)$. It supplies no proof authority.

\begin{proposition} \label{prop:standard-coordinate-dual-identification} \hlt{Standard Coordinate Dual Identification}\\
For $m\geq1$, define $\Theta:\mathbb{F}^m\to(\mathbb{F}^m)'$ by $(\forall\,\mathbf{y},\mathbf{z}\in\mathbb{F}^m)\,\Theta(\mathbf{y})(\mathbf{z})\defeq\mathbf{y}^t\mathbf{z}$. Then:
\begin{enumerate}[label=(\roman*)]
\setlength{\itemsep}{0pt}
\item $\Theta$ is an isomorphism;
\item $(\forall\,\psi\in(\mathbb{F}^m)')(\exists!\,\mathbf{y}\in\mathbb{F}^m)(\forall\,\mathbf{z}\in\mathbb{F}^m)\,\psi(\mathbf{z})=\mathbf{y}^t\mathbf{z}$.
\end{enumerate}
\begin{proof}
\begin{enumerate}[label=(\roman*)]
\setlength{\itemsep}{0pt}
\item For $a,b\in\mathbb{F}$ and $\mathbf{x},\mathbf{y},\mathbf{z},\mathbf{w}\in\mathbb{F}^m$, the identity $\Theta(\mathbf{y})(a\mathbf{z}+b\mathbf{w})=a\Theta(\mathbf{y})(\mathbf{z})+b\Theta(\mathbf{y})(\mathbf{w})$ holds by Proposition \ref{prop:matrix-product-associativity-and-identities}, so $\Theta(\mathbf{y})\in(\mathbb{F}^m)'$.\\
The identity $\Theta(a\mathbf{x}+b\mathbf{y})(\mathbf{z})=a\Theta(\mathbf{x})(\mathbf{z})+b\Theta(\mathbf{y})(\mathbf{z})$ holds by Proposition \ref{prop:elementary-transpose-rules} and Proposition \ref{prop:matrix-product-associativity-and-identities}; hence $\Theta$ is linear by Assumption \ref{assump:function-extensionality}. If $\Theta(\mathbf{y})=0$, then $(\forall\,j\in\{1,\ldots,m\})\,0=\Theta(\mathbf{y})(\mathbf{e}_j)=y_j$, so $\mathbf{y}=\mathbf{0}$ and $\Theta$ is injective by Proposition \ref{prop:injective-iff-kernel-zero}.\\
Since $(1)$ is a basis of $\mathbb{F}$ over itself, $\dim(\mathbb{F}^m)=m=\dim((\mathbb{F}^m)')$ by Theorem \ref{thm:dimension-of-product} and Theorem \ref{thm:dual-basis-and-dimension}(ii).\\
Thus $\Theta$ is surjective by Corollary \ref{cor:equal-dimension-kernel-range}(ii) and is an isomorphism by Proposition \ref{prop:linear-map-invertible-iff-bijective}.
\item Every $\psi\in(\mathbb{F}^m)'$ has a unique preimage $\mathbf{y}\in\mathbb{F}^m$ under $\Theta$ by (i), and $\psi(\mathbf{z})=\Theta(\mathbf{y})(\mathbf{z})=\mathbf{y}^t\mathbf{z}$ by the defining formula.
\end{enumerate}
\hlt{Structural Commentary}: Matrix identities prove linearity; coordinate evaluations and equal dimensions prove bijectivity.\\
\hlt{Key Insight}: Every coordinate functional is multiplication by one transposed column.
\end{proof}
\end{proposition}

\begin{corollary} \label{cor:matrix-equation-left-nullspace-alternative} \hlt{Matrix Equation Alternative}\\
Let $m,n\geq1$, $A\in\operatorname{Mat}(m,n,\mathbb{F})$, and $\mathbf{b}\in\mathbb{F}^m$. Exactly one holds:
\begin{enumerate}[label=(\roman*)]
\setlength{\itemsep}{0pt}
\item $(\exists\,\mathbf{x}\in\mathbb{F}^n)\,A\mathbf{x}=\mathbf{b}$;
\item $(\exists\,\mathbf{y}\in\mathbb{F}^m)[\mathbf{y}^tA=\mathbf{0}^t\land\mathbf{y}^t\mathbf{b}\neq0]$.
\end{enumerate}
\begin{proof}
Let $T_A:\mathbb{F}^n\to\mathbb{F}^m$ be the matrix-induced map.\\
Then $T_A\in\LL(\mathbb{F}^n,\mathbb{F}^m)$ by Proposition \ref{prop:matrix-induced-map-linear}, and both coordinate spaces are finite-dimensional by Theorem \ref{thm:dimension-of-product}.
\begin{enumerate}[label=(\roman*)]
\setlength{\itemsep}{0pt}
\item The equivalence $\mathbf{b}\in\operatorname{range}(T_A)\iff(\exists\,\mathbf{x}\in\mathbb{F}^n)\,A\mathbf{x}=\mathbf{b}$ holds by $T_A\mathbf{x}=A\mathbf{x}$.
\item By Proposition \ref{prop:standard-coordinate-dual-identification}(ii), each $\psi\in(\mathbb{F}^m)'$ has the unique form $\psi(\mathbf{z})=\mathbf{y}^t\mathbf{z}$. For this $\mathbf{y}$, $\psi T_A=0$ iff $(\forall\,\mathbf{x}\in\mathbb{F}^n)\,\mathbf{y}^tA\mathbf{x}=0$ iff $\mathbf{y}^tA=\mathbf{0}^t$ by evaluation on the standard basis of $\mathbb{F}^n$; also, $\psi(\mathbf{b})\neq0$ iff $\mathbf{y}^t\mathbf{b}\neq0$. Exactly one holds by Theorem \ref{thm:range-membership-dual-alternative}.
\end{enumerate}
\hlt{Structural Commentary}: Coordinate identification converts the abstract separating functional into a left-null vector.\\
\hlt{Key Insight}: Left-null vectors are coordinate dual certificates.
\end{proof}
\end{corollary}

\begin{theorem} \label{thm:matrix-of-dual-map-is-transpose} \hlt{Matrix of the Dual Map Is the Transpose}\\
Let $\dim(\mathcal{V})=n\geq1$ and $\dim(\mathcal{W})=m\geq1$, let $T\in\LL(\mathcal{V},\mathcal{W})$, and let $\beta=(\mathbf{v}_1,\ldots,\mathbf{v}_n)$ and $\gamma=(\mathbf{w}_1,\ldots,\mathbf{w}_m)$ be bases.\\
If $\beta'$ and $\gamma'$ are the corresponding dual bases, then $[T']_{\beta'\leftarrow\gamma'}=([T]_{\gamma\leftarrow\beta})^t$.
\begin{proof}
Set $A\defeq[T]_{\gamma\leftarrow\beta}$ and $C\defeq[T']_{\beta'\leftarrow\gamma'}$. Write $\beta'=(\varphi_1,\ldots,\varphi_n)$ and $\gamma'=(\psi_1,\ldots,\psi_m)$. By Definition \ref{def:matrix-representation-of-linear-map} and Definition \ref{def:dual-basis}, $(\forall\,k\in\{1,\ldots,n\})(\forall\,j\in\{1,\ldots,m\})\,C_{k,j}=(T'\psi_j)(\mathbf{v}_k)=\psi_j(T\mathbf{v}_k)=A_{j,k}$, where $T\mathbf{v}_k=\sum\limits_{r=1}^{m}A_{r,k}\mathbf{w}_r$. Hence $C=A^t$.\\
\hlt{Structural Commentary}: Evaluating dual-basis expansions identifies each dual-matrix entry with the transposed primal entry.\\
\hlt{Key Insight}: Dual evaluation swaps the row and column indices.
\end{proof}
\end{theorem}
